% concatenated from AGF_introduction_V92.tex
\documentclass[11pt,reqno]{amsart}
\usepackage[T1]{fontenc}
\usepackage[utf8]{inputenc}
\usepackage[english]{babel}
\usepackage[margin=1in, top=1in, bottom=1in]{geometry}
\usepackage{amsmath, amssymb, amsthm, mathtools}
\usepackage{tikz-cd}
\usepackage{enumitem}
\usepackage{hyperref}

\newtheorem{theorem}{Theorem}[section]
\newtheorem{lemma}[theorem]{Lemma}
\newtheorem{proposition}[theorem]{Proposition}
\newtheorem{corollary}[theorem]{Corollary}
\theoremstyle{definition}
\newtheorem{definition}[theorem]{Definition}
\newtheorem{remark}[theorem]{Remark}
\newtheorem{example}[theorem]{Example}
\newtheorem{construction}[theorem]{Construction}
\newtheorem{assumption}[theorem]{Assumption}
\newtheorem{problem}{Open Problem}
\newtheorem{principle}{Principle}[section]

\newcommand{\N}{\mathbb{Z}_{\ge 0}}
\newcommand{\Z}{\mathbb{Z}}
\newcommand{\Q}{\mathbb{Q}}
\newcommand{\C}{\mathbb{C}}
\newcommand{\R}{\mathbb{R}}
\newcommand{\dd}{\mathrm{d}}
\newcommand{\ee}{\mathrm{e}}
\DeclareMathOperator{\ord}{ord}

\hypersetup{colorlinks=true, linkcolor=blue, citecolor=blue, urlcolor=blue}
\numberwithin{equation}{section}

\title[The holonomic triangle]{The holonomic triangle: from a symmetry between $\ee$ and $\pi$ to additive Gamma functions}
\author{Benoit Cloitre}
\date{\today}

\begin{document}
\maketitle

\begin{abstract}
Two linear recurrences exhibit mirror symmetry connecting the constants $\ee$ and $\pi$. When parametrized, their asymptotic connection constants extend to meromorphic functions satisfying additive functional equations with rational coefficients. We call such functions \emph{additive Gamma functions} (AGFs), recognizing Euler's $\Gamma(z)$ as the order-1 prototype. Our theory reveals a structural dichotomy: one AGF is expressible as Gamma ratios (regular case), another involves incomplete Gamma (irregular case).  AGFs complete a holonomic triangle between P-recursive sequences, additive functional equations, and differential equations, unifying discrete and continuous perspectives under the condition that Gamma factors in asymptotics have integer slopes.
\end{abstract}

\subjclass[2020]{Primary: 11B37, 33B15; Secondary: 39B32, 30E15}
\keywords{Additive functional equations, P-recursive sequences, Connection constants, Gamma function, Holonomic systems}

\section{Introduction}

\subsection*{Notation}

Throughout this paper, $\N := \{0,1,2,\ldots\}$ denotes the nonnegative integers, and $\Z$, $\Q$, $\R$, $\C$ denote the standard number systems. We use the principal branch for $\log$ and $\Gamma$, with the cut along $(-\infty,0]$ for $\log$ and poles at $\{0,-1,-2,\ldots\}$ for $\Gamma$.

We use $n$ for discrete indices and $m\in\N$ for the discrete parameter in the recurrences. Analytic continuation is always performed in a complex variable denoted by $z$ (or occasionally $w$), so that no argument of a meromorphic function is ever written with the symbol $n$ unless it is explicitly an integer.


\subsection{The mirror symmetry}

This work originates from a mirror symmetry observed by the author in 2002 between two linear recurrences connecting the mathematical constants $\ee$ and $\pi$:
\begin{equation}\label{eq:mirror}
\boxed{
\begin{aligned}
u_{n+2}&=u_{n+1}+\frac{u_{n}}{n}, \qquad &\text{\small (world of $\ee$)}\\[0.3em]
v_{n+2}&=\frac{v_{n+1}}{n}+v_{n}, \qquad &\text{\small (world of $\pi$)}
\end{aligned}
}
\end{equation}
both with initial conditions $u_{1}=v_{1}=0$ and $u_{2}=v_{2}=1$ for $n\ge 1$. These recurrences satisfy
$$\lim_{n\to\infty} \frac{n}{u_n} = \ee, \qquad \lim_{n\to\infty} \frac{2n}{v_n^2} = \pi.$$

The sequence $(u_n)$ relates to derangement numbers \cite{OEIS_A000166} while $(v_n)$ connects to Wallis integrals and double factorials \cite{OEIS_A006882}; we refer to Loya \cite[Section 7.7.3]{Loya2017} for detailed derivations. Finch \cite{Finch2003} notes these recurrences as a counterpoint to Wimp's dichotomy \cite{Wimp1984}, which distinguishes ``$\ee$-mathematics'' (linear, explicit) from ``$\pi$-mathematics'' (nonlinear, mysterious).

While each sequence individually belongs to mathematical folklore, their juxtaposition suggests a deeper structure connecting $\ee$ and $\pi$. The following arithmetic duality begins to reveal this hidden structure.

\subsection{The arithmetic duality}

By introducing a parameter $m\in\N$ into the denominators of \eqref{eq:mirror}:
\begin{equation}\label{eq:param}
\boxed{
\begin{aligned}
u_{n+2}(m) &= u_{n+1}(m)+\frac{u_{n}(m)}{n+m}, \\[0.3em]
v_{n+2}(m) &= \frac{v_{n+1}(m)}{n+m}+v_{n}(m),
\end{aligned}
}
\end{equation}
we obtain parametrized sequences $u_n(m)$ and $v_n(m)$ with asymptotic behaviors
\[
u_n(m) \sim f(m) \cdot n, \qquad v_n(m) \sim g(m) \cdot \sqrt{n}.
\]

The analysis of these connection constants for $m\in\N$ reveals an arithmetic duality. Once normalized by $f(0)=\frac{1}{\ee}$ and $g(0)=\sqrt{\frac{2}{\pi}}$, they take the form of linear combinations:
\begin{equation}\label{eq:duality}
\boxed{
\begin{aligned}
(-1)^m \frac{f(m)}{f(0)} &= a_m - \ee \cdot b_m, \quad \text{with } a_m, b_m \in \Z, \\[0.5em]
(-1)^m \frac{g(m)}{g(0)} &= p_m - \pi \cdot q_m, \quad \text{with } p_m, q_m \in \Q.
\end{aligned}
}
\end{equation}

The explicit formulas for $a_m$, $b_m$, $p_m$, $q_m$ are established in Theorems~\ref{thm:e-world} and~\ref{thm:pi-world}.

This arithmetic duality suggests that $f(m)$ and $g(m)$ are not merely real sequences but discrete samplings of underlying analytic functions. The question naturally arises: can these connection constants be extended to the complex plane, and if so, do the resulting functions inherit algebraic structure from the recurrences?

\subsection{From arithmetic linear forms to meromorphic functions}

From this point on, we keep the letters $m,n$ for integers (parameters or discrete indices) and reserve $z$ for a complex variable.

When analytically continued to the complex plane, these constants extend to meromorphic functions $f(z)$ and $g(z)$ satisfying additive 
functional equations with rational coefficients. This motivates the central definition of this work:

\begin{definition}[Additive Gamma function]\label{def:AGF-intro}
A meromorphic function $h:\C\to\C$ is an additive Gamma function (AGF) of order $r$ if it satisfies the following three conditions:
\begin{enumerate}[label=(\roman*)]
\item \textbf{Additive functional equation:} $h$ satisfies a homogeneous additive functional equation of order $r$ with rational coefficients:
$$\sum_{k=0}^r R_k(z) h(z+k) = 0,$$
where $R_k(z) \in \Q(z)$ are rational functions with $R_0 \not\equiv 0$ and $R_r \not\equiv 0$, ensuring the equation has effective order~$r$.

\item \textbf{Moderate growth:} For any vertical strip $a \leq \mathrm{Re}\, z \leq b$ not containing poles of $h$, there exist constants $C_{a,b} > 0$ and $N_{a,b} \in \N$ such that
$$|h(z)| \le C_{a,b}(1+|\mathrm{Im}\, z|)^{N_{a,b}}$$
for all $z$ in the strip.

\item \textbf{Normalization:} For any choice of $z_0 \in \R$ avoiding poles of $h$, the function $h$ is uniquely determined by the functional equation (i), the growth condition (ii), and its values on the real segment $[z_0, z_0+r)$.
\end{enumerate}
\end{definition}

\begin{remark}
Condition (iii) ensures uniqueness: if two functions $h_1$ and $h_2$ satisfy the same additive functional equation (i) with moderate growth (ii) and agree on a real segment $[z_0, z_0+r)$, then $h_1 \equiv h_2$ by Nörlund's uniqueness theorem \cite[Chapter IV, Theorem 1]{Norlund1924}. This prevents ambiguities from periodic functions of period~$r$. In particular, if the growth hypothesis (ii) is dropped, one can generally add to a given solution a nontrivial $r$-periodic meromorphic function (or, more broadly, a periodic solution of the same shift-recurrence), so the normalization (iii) alone would no longer single out a unique AGF. For instance, $f(z)$ (an order-2 AGF) is uniquely determined by its values $f(0)$ and $f(1)$ on the interval $[0,2)$, together with the AFE of Theorem~\ref{thm:e-world}.
\end{remark}

The functions $f(z)$ and $g(z)$ introduced above are order-2 AGFs, while Euler's $\Gamma(z)$ (satisfying $\Gamma(z+1) = z\Gamma(z)$) serves as the order-1 prototype.

\subsection{Main results}

The heart of the paper is the following characterization, proved in Section~\ref{sec:general-theory} 
(Theorem~\ref{thm:AGF-equivalence}): a meromorphic function is an additive Gamma 
function of order $r$ if and only if it arises as a connection constant from a 
sequence holonomic in $(n,z)$ with integer-slope asymptotics:
$$u_n(z) \sim \Lambda(n,z)h(z) \quad \text{as } n \to \infty,$$
where $\Lambda(n,z) = \lambda^n n^{\rho(z)} \prod_j \Gamma(n+\alpha_j z+\beta_j)^{m_j}$
with constant $\lambda\in\C^*$, affine $\rho(z)$ with integer slope, integer slopes $\alpha_j\in\Z$, and $m_j\in\Z$.

\paragraph{Uniqueness and the regular/irregular dichotomy.}
These functional equations are uniquely determined by their values on a compact interval of length 2 together with the growth conditions.
By Nörlund's theorem (Lemma~\ref{lem:unicity-norlund}),
this uniqueness holds under our vertical-strip growth assumptions.
The explicit formulas reveal a structural dichotomy: $f(z)$ involves the incomplete Gamma function (irregular case), while $g(z)$ is expressed as Gamma ratios (regular case).

\subsubsection*{Motivating examples: the worlds of $\ee$ and $\pi$}

The first two theorems below instantiate the general mechanism in order $r=2$. 
They are meant to be read as \emph{illustrations} of the general theory: they exhibit the two ``archetypal'' AGFs that motivate the
definitions (incomplete Gamma vs.\ Gamma ratios) and foreshadow the regular/irregular dichotomy analyzed later in 
Section~\ref{sec:dichotomy}.

\begin{theorem}[Connection constant in the world of $\ee$]\label{thm:e-world}
For the parametrized sequence $u_n(m)$ defined by \eqref{eq:param}, there exists a connection constant $f(m)$ such that 
\[
u_n(m) \sim f(m) \cdot n \quad \text{as } n\to\infty.
\]
The function $f(m)$ extends to a meromorphic function $f(z)$ on $\C$ satisfying the additive functional equation
\begin{equation}\label{eq:f-AFE-thm}
f(z+2) = (z+2)f(z) - (z+2)f(z+1).
\end{equation}
The function $f(z)$ admits an explicit formula in terms of the incomplete Gamma function $\gamma(a,x)$ and the confluent hypergeometric function ${}_1F_1$ (see \eqref{eq:f-explicit} and \eqref{eq:f-confluent} in Section~\ref{sec:e-world}), exhibits polynomial growth in vertical strips, and extends meromorphically to $\C$ with poles at $z \in \{-2,-3,-4,\ldots\}$.

The arithmetic structure \eqref{eq:duality} is realized with $a_m = (m+1)!$ and $b_m = D_{m+1}$, the derangement numbers.
\end{theorem}

The proof is given in Section~\ref{sec:e-world}.

\begin{theorem}[Connection constant in the world of $\pi$]\label{thm:pi-world}
For the parametrized sequence $v_n(m)$ defined by \eqref{eq:param}, there exists a connection constant $g(m)$ such that
\[
v_n(m) \sim g(m) \cdot \sqrt{n} \quad \text{as } n\to\infty.
\]
The function $g(m)$ extends to a meromorphic function $g(z)$ on $\C$ satisfying the additive functional equation
\begin{equation}\label{eq:g-AFE-thm}
g(z+2) = g(z) - \frac{g(z+1)}{z+1}.
\end{equation}
The function $g(z)$ admits an explicit formula as a combination of Gamma ratios (see \eqref{eq:g-explicit} in Section~\ref{sec:pi-world}), exhibits polynomial growth in vertical strips, and extends meromorphically to $\C$ with simple poles at the negative integers $z \in \{-1,-2,-3,\ldots\}$ and is holomorphic at $z=0$.

The arithmetic structure \eqref{eq:duality} is realized with explicit formulas involving double factorials (see Proposition~\ref{prop:pq-formulas}).
\end{theorem}

The proof is given in Section~\ref{sec:pi-world}.

\subsection{Related work and positioning}

This work connects classical asymptotic analysis, modern holonomic systems theory, and analytic combinatorics through additive functional equations.

\textbf{Classical foundations.} Nörlund \cite{Norlund1924} developed difference calculus and uniqueness theorems we invoke, but did not connect to P-recursive asymptotics. Birkhoff-Trjitzinsky \cite{BirkhoffTrjitzinsky1933} analyzed asymptotic expansions of difference equations comprehensively, but treated connection constants as byproducts rather than primary objects of study. Barnes \cite{Barnes1901} extended Gamma via \emph{multiplicative} equations; AGFs provide an \emph{additive} counterpart.

\textbf{Holonomic systems.} Zeilberger \cite{Zeilberger1990,Zeilberger1991} and Chyzak-Salvy \cite{ChyzakSalvy1998} developed algorithmic approaches where connection constants appear as byproducts. Our Integer Slope Condition makes explicit an arithmetic constraint implicit in their elimination algorithms. Kauers and Paule \cite{KauersPaule2011} provide a comprehensive modern treatment of these computational techniques.

\textbf{Analytic combinatorics.} Flajolet-Sedgewick \cite{FlajoletSedgewick2009} provides the singularity analysis underlying our constructions. Stanley \cite{Stanley1999} treats the combinatorial sequences appearing in our examples.

\subsection{Organization}

The paper is organized as follows: 

Section~\ref{sec:mirror} establishes technical preliminaries including Carlson's theorem and uniqueness results.

Section~\ref{sec:e-world} develops the asymptotic analysis for the world of $\ee$, deriving the generating function, extracting the connection constant $f(m)$, establishing the discrete recurrence and arithmetic structure, performing analytic continuation, and proving the additive functional equation. This establishes Theorem~\ref{thm:e-world}.

Section~\ref{sec:pi-world} develops the parallel asymptotic analysis for the world of $\pi$, following the same program adapted to $\sqrt{n}$ asymptotics and $\Q$-linear forms. This establishes Theorem~\ref{thm:pi-world}.

Section~\ref{sec:gamma-prototype} verifies that $\Gamma(z)$, $f(z)$, and $g(z)$ satisfy the three defining properties of additive Gamma functions, establishing the growth bounds invoked in Sections~\ref{sec:e-world} and~\ref{sec:pi-world}. This section reveals the AGF hierarchy with Euler's $\Gamma(z)$ as the order-1 prototype.

Section~\ref{sec:general-theory} develops the general theory of AGFs, proving the Integer Slope Condition (Lemma~\ref{lem:integer-slope}) and the characterization theorem (Theorem~\ref{thm:AGF-equivalence}).

Section~\ref{sec:dichotomy} analyzes the regular-irregular dichotomy between $f(z)$ (confluent hypergeometric) and $g(z)$ (ordinary hypergeometric).

Section~\ref{sec:holonomic} presents the holonomic triangle (Figure~\ref{fig:triangle}, Principle~\ref{prin:holonomic}), providing a conceptual framework connecting P-recurrences, AFEs, and D-finite ODEs under the integer slope condition.

We conclude with open questions and future directions.

\section{Technical preliminaries}\label{sec:mirror}

The method employed in this paper to derive additive functional equations from discrete recurrences relies on Carlson's theorem, which we state for reference:

\begin{theorem}[Carlson]\label{thm:carlson}
Let $\varphi$ be a holomorphic function on the half-plane $\Re z \geq 0$ such that:
\begin{enumerate}[label=(\roman*)]
\item $\varphi(n) = 0$ for all integers $n \in \N$,
\item $|\varphi(z)| \leq C \ee^{\tau |\mathrm{Im}\, z|}$ for some constants $C > 0$ and $\tau < \pi$,
\end{enumerate}
Then $\varphi \equiv 0$ throughout the half-plane $\Re z \geq 0$. \textup{(See \cite[Theorem 6.20]{Titchmarsh1939}.)}
\end{theorem}

This theorem provides a powerful tool for proving functional equations: if a function $H(z)$ satisfies a discrete recurrence $H(m) = 0$ for all $m \in \N$ and has controlled exponential growth, then $H \equiv 0$, establishing the desired functional equation.

Uniqueness for solutions of additive functional equations follows from Nörlund's classical theorem:

\begin{lemma}[Uniqueness under AFE]\label{lem:unicity-norlund}
Let $\phi_1,\phi_2$ be meromorphic on $\C$, of exponential type $<\pi$ on every
vertical strip, satisfying the same additive functional equation
\[
  \mathcal{L}(T)\phi(z) := \sum_{j=0}^{r} a_j(z)\,\phi(z+j) \;=\; 0
\]
with polynomial (or rational) coefficients $a_j$, and coinciding on a real
interval of positive length. Then $\phi_1\equiv \phi_2$.
\end{lemma}

\begin{proof}
By Nörlund's uniqueness theorem for difference equations
(\emph{Vorlesungen über Differenzenrechnung} \cite{Norlund1924}, Ch.~IV, Th.~1),
the exponential-type bound $<\pi$ together with equality on a full real
interval forces $\phi_1\equiv \phi_2$.
\end{proof}

These tools will be applied in Sections~\ref{sec:e-world} and~\ref{sec:pi-world} to establish functional equations, with the required growth bounds verified in Section~\ref{sec:gamma-prototype}.

\section{Asymptotic analysis: the world of $\ee$}\label{sec:e-world}

We now establish Theorem~\ref{thm:e-world} through detailed asymptotic analysis, deriving the generating function, extracting the connection constant, and proving the additive functional equation.
The holonomic differential equation satisfied by $U_m$ is recorded in Proposition~\ref{prop:GF-ODE-e}.

\subsection{Generating function and differential equation}

Let $U_m(x) = \sum_{n=1}^\infty u_n(m) x^n$. Multiplying the recurrence $u_{n+2}(m) = u_{n+1}(m) + \frac{u_n(m)}{n+m}$ by $x^n$ and summing over $n \ge 1$ yields
\[
\sum_{n=1}^\infty u_{n+2}(m)x^n = \sum_{n=1}^\infty u_{n+1}(m)x^n + \sum_{n=1}^\infty \frac{u_n(m)}{n+m}x^n.
\]
With the initial conditions $u_1(m)=0$ and $u_2(m)=1$, the left side equals $\frac{U_m(x) - u_1x - u_2x^2}{x^2} = \frac{U_m(x) - x^2}{x^2}$, and the first term on the right equals $\frac{U_m(x) - u_1x}{x} = \frac{U_m(x)}{x}$. Setting $J_m(x) = \sum_{n=1}^\infty \frac{u_n(m)}{n+m}x^n$ and using $\int_0^x t^{n+m-1}\, \dd t = \frac{x^{n+m}}{n+m}$:
\[
J_m(x) = x^{-m}\int_0^x U_m(t)t^{m-1}\, \dd t.
\]
Substituting into the equation $\frac{U_m(x) - x^2}{x^2} = \frac{U_m(x)}{x} + J_m(x)$ and rearranging yields the integro-differential relation
\[
(1-x)U_m(x) = x^2(1 + J_m(x)).
\]

Differentiating this equation with respect to $x$ and eliminating $J_m(x)$ leads to a first-order linear ODE:
\[
(1-x)U'_m(x) = U_m(x) \left(\frac{x^2-x+2-m+mx}{x}\right) + mx.
\]

Via standard methods, this ODE has integrating factor $\mu(x) = x^{m-2}(1-x)^2\ee^x$, yielding:
\[
U_m(x) = \frac{x^{2-m}}{(1-x)^2\ee^x}\left(\delta_{m,0} + \int_0^x mt^{m-1}(1-t)\ee^t\, \dd t\right).
\]

\begin{remark}[Order of the ODE]\label{rem:order-ode}
Although the P-recurrence for $u_n(m)$ has width (order) $r=2$, the resulting ODE for the ordinary generating function $U_m(x)$ has order 1. When the recurrence is written in standard P-recursive form (by clearing denominators):
$$(n+m)u_{n+2}(m) - (n+m)u_{n+1}(m) - u_n(m) = 0,$$
the polynomial coefficients have maximum degree $d=1$. 

More generally, for a P-recursive sequence of width $r$ with coefficient polynomials of maximum degree $d$, under suitable growth conditions ensuring the ordinary generating function (OGF) is D-finite (which holds for sequences with moderate algebraic growth, as is the case for AGFs with integer-slope asymptotics), the OGF satisfies a linear ODE whose order depends on both $r$ and $d$. Based on elimination theory in Ore algebras \cite{ChyzakSalvy1998,Koutschan2010}, the order is typically at most $(d+1)r$ in the worst case, though often much lower in practice. For comparison, the exponential generating function (EGF) satisfies an ODE of order at most $r$ (Stanley \cite{Stanley1999}, Ch.~6). This phenomenon is a fundamental feature of the holonomic triangle (see Principle~\ref{prin:holonomic} and Figure~\ref{fig:triangle}).
\end{remark}

\subsection{Singularity analysis and asymptotic behavior}

As $x\to 1^-$: 
$$U_m(x) \sim \frac{J_m}{\ee(1-x)^{2}} \quad\text{where}\quad J_m = \delta_{m,0} + \int_0^1 mt^{m-1}(1-t)\ee^t\, \dd t.$$

The point $x=1$ is a pole of order 2 (dominant singularity). By the transfer theorem of Flajolet and Sedgewick \cite{FlajoletSedgewick2009} (Theorem VI.3, p.393), if a generating function has a singular expansion $G(x) \sim C(1-x)^{-\alpha}$ near $x=1$ with $\alpha = 2$, then
$$[x^n]G(x) \sim \frac{C}{\Gamma(\alpha)} n^{\alpha-1} = C \cdot n.$$

Applying this with $C = \ee^{-1}J_m$ yields
$$u_n(m) \sim \ee^{-1}J_m \cdot n.$$

This establishes the asymptotic behavior claimed in Theorem \ref{thm:e-world} with the connection constant
\begin{equation}\label{eq:f-def}
f(m) = \ee^{-1}J_m.
\end{equation}

\subsection{Discrete recurrence and arithmetic structure}

We derive the recurrence satisfied by $J_m = f(m)\ee$. From the definition above, we have
\begin{equation}\label{eq:J-def}
J_m = \begin{cases}
1 & \text{if } m=0,\\
\displaystyle\int_0^1 mt^{m-1}(1-t)\ee^t\, \dd t & \text{if } m \geq 1.
\end{cases}
\end{equation}

For $m \geq 1$, we analyze the integral expression using a key identity.

\textit{Auxiliary integrals.} Define the auxiliary integrals $I_m = \int_0^1 t^m\ee^t\, \dd t$ for $m \geq 0$. We have $I_0 = \ee-1$. Integration by parts ($u=t^m, \dd v=\ee^t\, \dd t$) yields the well-known recurrence:
\begin{equation}\label{eq:I-rec}
I_m = [t^m\ee^t]_0^1 - m\int_0^1 t^{m-1}\ee^t\, \dd t = \ee - mI_{m-1}, \quad m \geq 1.
\end{equation}

\textit{Relating $J_m$ and $I_m$.} We can express $J_m$ (for $m \geq 1$) directly in terms of $I_m$:
\begin{align*}
J_m &= \int_0^1 mt^{m-1}(1-t)\ee^t\, \dd t = m\int_0^1 (t^{m-1}\ee^t - t^m\ee^t)\, \dd t \\
&= m(I_{m-1} - I_m).
\end{align*}

We now establish a crucial identity. Using the recurrence \eqref{eq:I-rec} for $I_m$:
\begin{align*}
J_m &= m(I_{m-1} - (\ee - mI_{m-1})) = m((m+1)I_{m-1} - \ee).
\end{align*}

Now consider $I_{m+1}$. Using the recurrence \eqref{eq:I-rec} repeatedly:
\begin{align*}
I_{m+1} &= \ee - (m+1)I_m \\
&= \ee - (m+1)(\ee - mI_{m-1}) \\
&= m(m+1)I_{m-1} - m\ee.
\end{align*}

Thus, we have proven the identity:
\begin{equation}\label{eq:J-I-identity}
J_m = I_{m+1} \quad \text{for } m \geq 1.
\end{equation}

We verify this for $m=0$: $J_0=1$ (by definition), and $I_1 = \ee - 1\cdot I_0 = \ee - (\ee-1) = 1$. The identity holds for all $m \geq 0$.

\textit{Deriving the recurrence for $J_m$.} Using the identity \eqref{eq:J-I-identity} and the recurrence \eqref{eq:I-rec}:
\begin{align*}
J_{m+1} &= I_{m+2} = \ee - (m+2)I_{m+1} = \ee - (m+2)J_m.
\end{align*}

This establishes the first-order recurrence:
\begin{equation}\label{eq:J-rec}
J_{m+1} = \ee - (m+2)J_m \quad \text{for } m \geq 0.
\end{equation}

\textit{Elimination to second order.} To obtain a homogeneous recurrence of order 2, we eliminate the inhomogeneous term $\ee$ by taking successive instances of \eqref{eq:J-rec}:
$$J_{m+2} = \ee - (m+3)J_{m+1}, \quad J_{m+1} = \ee - (m+2)J_m.$$

Subtracting these equations:

$$J_{m+2} = (m+2)J_m - (m+3)J_{m+1} + J_{m+1} = (m+2)[J_m - J_{m+1}].$$

Since $f(m) = \ee^{-1}J_m$, dividing by $\ee$ gives:
\begin{equation}\label{eq:f-discrete}
f(m+2) = (m+2)[f(m) - f(m+1)].
\end{equation}

\textit{Arithmetic structure.} We now derive the $\Z$-linear form in $\ee$ by introducing the sign-normalized quantity:
$$K_m = (-1)^m J_m.$$

From the first-order recurrence $J_{m+1} = \ee - (m+2)J_m$, multiplying by $(-1)^{m+1}$:
$$(-1)^{m+1}J_{m+1} = (-1)^{m+1}\ee - (-1)^{m+1}(m+2)J_m = -(-1)^m\ee + (m+2)(-1)^mJ_m.$$

In terms of $K_m$, this becomes:
$$K_{m+1} = (m+2)K_m - (-1)^m \ee.$$

We seek a decomposition $K_m = a_m - \ee b_m$ with $a_m, b_m \in \Z$ and $a_m, b_m \ge 0$. Substituting into the recurrence:
$$a_{m+1} - \ee b_{m+1} = (m+2)(a_m - \ee b_m) - (-1)^m \ee.$$

Expanding and separating terms by the $\Q$-linear independence of $\{1, \ee\}$ (using the transcendence of $\ee$):
$$\begin{cases}
a_{m+1} = (m+2)a_m, \\
b_{m+1} = (m+2)b_m + (-1)^m.
\end{cases}$$

With initial conditions $K_0 = 1$ (giving $a_0=1, b_0=0$) and $K_1 = 2-\ee$ (giving $a_1=2, b_1=1$), we find:
$$a_m = (m+1)!, \quad b_m = D_{m+1},$$
where $D_m$ are the derangement numbers. This confirms the $\Z$-linear form in \eqref{eq:duality}.

\subsection{Analytic continuation and functional equation}

We define the analytic continuation via the auxiliary function $I(z) = \int_0^1 t^z\ee^t\, \dd t$, which converges and is holomorphic for $\Re z > -1$. By the identity principle for holomorphic functions, the relation $J_m = I_{m+1}$ established in \eqref{eq:J-I-identity} extends to the complex domain: $J(z) = I(z+1)$ for $\Re z > -1$. The connection constant is $f(z) = \ee^{-1}J(z)$. This provides an analytic continuation of $f(z)$ to the half-plane $\Re z > -2$.

\begin{lemma}[Growth bounds for $f$]\label{lem:growth-f}
The function $f(z)$ satisfies $|f(z)| = O(1/|z|)$ for $\Re z \geq 0$, ensuring the moderate growth required for Carlson's theorem. The proof, based on the explicit formula \eqref{eq:f-explicit} established below, is given in Section~\ref{sec:gamma-prototype}.
\end{lemma}

To derive the functional equation, we apply Carlson's theorem (Theorem~\ref{thm:carlson}). From the discrete recurrence \eqref{eq:f-discrete}, define 
$$h(z) = f(z+2) - (z+2)[f(z) - f(z+1)].$$

By \eqref{eq:f-discrete}, we have $h(m)=0$ for all $m\in\N$. By Lemma \ref{lem:growth-f}, $f(z)$ is holomorphic for $\Re z \geq 0$ and satisfies $|f(z)| = O(1/|z|)$. From this asymptotic expansion $f(z) = 1/(z+2) + O(1/z^2)$ as $|z|\to\infty$, we obtain $f(z)-f(z+1) = O(1/|z|^2)$. Therefore $(z+2)[f(z)-f(z+1)] = O(1/|z|)$, so that $h(z) = O(1/|z|)$ is bounded for $\Re z \geq 0$.

Since $h(z)$ is bounded, it has exponential type $0 < \pi$, so Carlson's theorem applies directly. Together with $h(n)=0$ for all $n\in\N$, we conclude $h(z) \equiv 0$ and thus:
\begin{equation}\label{eq:f-AFE}
f(z+2) = (z+2)f(z) - (z+2)f(z+1).
\end{equation}

\subsection{Explicit formula}

The incomplete Gamma function $\gamma(a,x) = \int_0^x t^{a-1}\ee^{-t}\, \dd t$ (see \cite[Section 8.2]{DLMF}, \cite[Section 8]{Temme1996}, or classically \cite[Section 6.5]{AbramowitzStegun}) is analytic in $x$ on the complex plane cut along $(-\infty, 0]$. Since $x=-1$ lies on this branch cut, we employ the standard analytic continuation and use the principal branch throughout. We now derive an explicit formula for $f(z)$ in terms of $\gamma$.

Recall that $f(z) = \ee^{-1}J(z) = \ee^{-1}I(z+1)$ where $I(z) = \int_0^1 t^z \ee^t\,\dd t$. The substitution $t = -u$ in the integral defining $\gamma(z+2, -1)$ gives:
\begin{align*}
\gamma(z+2, -1) &= \int_0^{-1} t^{z+1}\ee^{-t}\,\dd t \\
&= \int_0^1 (-u)^{z+1}\ee^u\,(-\, \dd u) \\
&= -\ee^{i\pi(z+1)} I(z+1),
\end{align*}
where we use the principal branch $(-1)^{z+1} = \ee^{i\pi(z+1)}$. Since $f(z) = \ee^{-1}I(z+1)$ and $\ee^{i\pi} = -1$, we obtain:
$$I(z+1) = -\ee^{-i\pi(z+1)}\gamma(z+2,-1) = -\ee^{-i\pi z} \cdot \ee^{-i\pi}\gamma(z+2,-1) = -\ee^{-i\pi z} \cdot (-1)\gamma(z+2,-1) = \ee^{-i\pi z}\gamma(z+2,-1).$$

Therefore:
\begin{equation}\label{eq:f-explicit}
f(z) = \ee^{-1-i\pi z}\gamma(z+2,-1).
\end{equation}

Alternatively, using the relation between incomplete Gamma and confluent hypergeometric functions \cite[Eq.~8.4.15]{DLMF}:
\begin{equation}\label{eq:f-confluent}
f(z) = \frac{1}{z+1}\,{}_1F_1(2; z+2; -1).
\end{equation}

The confluent hypergeometric function ${}_1F_1(a;b;x)$ satisfies the identity (Kummer's transformation):
$${}_1F_1(a;b;x) = \ee^x {}_1F_1(b-a;b;-x),$$
which can be verified to reproduce the functional equation \eqref{eq:f-AFE}.

For $z=0$: $\gamma(2,-1) = \int_0^{-1} u\ee^{-u}\, \dd u = 1$. Then $\ee^{-1-i\pi\cdot 0}\cdot(1) = \ee^{-1}$, hence $f(0) = \ee^{-1}$ as expected.

This completes the $\ee$-world analysis and establishes Theorem~\ref{thm:e-world}.

\section{Asymptotic analysis: the world of $\pi$}\label{sec:pi-world}

We now establish Theorem~\ref{thm:pi-world}, following the same analytical program as for the $\ee$-world. The argument parallels Section~\ref{sec:e-world} with adaptations for the $\sqrt{n}$ asymptotics and $\Q$-linear forms characteristic of the $\pi$-world. The holonomic differential equation satisfied by $V_m$ is recorded in Proposition~\ref{prop:GF-ODE-pi}.

\subsection{Generating function and differential equation}

For $(v_n(m))$, let $V_m(x) = \sum_{n=1}^\infty v_n(m)x^n$. Following analysis similar to Section~\ref{sec:e-world}, the first-order ODE is:
\[
(1-x^2)V_m'(x) + V_m(x)\left(\frac{m-2}{x}-mx-1\right) = mx.
\]

Solving via integrating factor $\mu(x) = x^{m-2}(1-x)^{3/2}(1+x)^{1/2}$ yields:
\[
V_m(x) = \frac{x^{2-m}}{(1-x)^{3/2}(1+x)^{1/2}}\left(\delta_{m,0} + \int_0^x mt^{m-1}\sqrt{\frac{1-t}{1+t}}\, \dd t\right).
\]

\subsection{Singularity analysis and asymptotic behavior}

As $x\to 1^-$: 
$$V_m(x) \sim \frac{L_m}{\sqrt{2}(1-x)^{3/2}} \quad\text{where}\quad L_m = \delta_{m,0} + \int_0^1 mt^{m-1}\sqrt{\frac{1-t}{1+t}}\, \dd t.$$

The point $x=1$ is the dominant singularity with exponent $\alpha = \frac{3}{2}$. By the transfer theorem of Flajolet and Sedgewick \cite{FlajoletSedgewick2009} (Theorem VI.1, p.388) with $\Gamma\left(\frac{3}{2}\right) = \frac{\sqrt{\pi}}{2}$:
$$v_n(m) \sim \frac{L_m/\sqrt{2}}{\Gamma(3/2)} n^{3/2-1} = \sqrt{\frac{2}{\pi}}L_m \sqrt{n}.$$

This establishes the asymptotic behavior claimed in Theorem \ref{thm:pi-world} with the connection constant
\begin{equation}\label{eq:g-def}
g(m) = \sqrt{\frac{2}{\pi}}L_m.
\end{equation}

\subsection{Discrete recurrence: derivation}

We derive the recurrence satisfied by $L_m = g(m)\sqrt{\pi/2}$. From the definition above, we have
\begin{equation}\label{eq:L-def}
L_m = \begin{cases}
1 & \text{if } m=0,\\
\displaystyle\int_0^1 mt^{m-1}\sqrt{\frac{1-t}{1+t}}\,\dd t & \text{if } m \geq 1.
\end{cases}
\end{equation}

For $m \geq 1$, we analyze the integral expression.

Define $w(t) = \sqrt{\frac{1-t}{1+t}}$. We compute the derivative of $w(t)$ by taking the logarithmic derivative:
$$\frac{w'(t)}{w(t)} = \frac{\dd}{\, \dd t} \left(\frac{1}{2}\ln(1-t) - \frac{1}{2}\ln(1+t)\right) = -\frac{1}{2(1-t)} - \frac{1}{2(1+t)} = -\frac{1}{1-t^2}.$$

Thus, $w'(t) = -\frac{w(t)}{1-t^2}$.

For $m \geq 1$, we apply integration by parts to $L_m = \int_0^1 mt^{m-1}w(t)\, \dd t$. Setting $u' = mt^{m-1}$ (so $u=t^m$) and $v = w(t)$, we obtain:
\begin{align*}
L_m &= [t^m w(t)]_0^1 - \int_0^1 t^m w'(t)\, \dd t.
\end{align*}

The boundary term is $[1\cdot 0 - 0\cdot 1] = 0$, hence:
\begin{align*}
L_m &= - \int_0^1 t^m \left(-\frac{w(t)}{1-t^2}\right)\, \dd t = \int_0^1 \frac{t^m w(t)}{1-t^2}\, \dd t.
\end{align*}

This identity is valid for $m \geq 1$. We now relate $L_m$ and $L_{m+2}$ using this identity:
\begin{align*}
L_m - L_{m+2} &= \int_0^1 \frac{t^m w(t)}{1-t^2}\, \dd t - \int_0^1 \frac{t^{m+2} w(t)}{1-t^2}\, \dd t \\
&= \int_0^1 \frac{w(t)}{1-t^2} (t^m - t^{m+2})\, \dd t \\
&= \int_0^1 t^m w(t)\, \dd t.
\end{align*}

By definition \eqref{eq:L-def}, the right-hand side equals $L_{m+1}/(m+1)$:
$$L_{m+1} = \int_0^1 (m+1)t^m w(t)\, \dd t.$$

Combining these results yields:
$$L_m - L_{m+2} = \frac{L_{m+1}}{m+1}.$$

Rearranging gives the homogeneous second-order recurrence:
\begin{equation}\label{eq:L-rec}
L_{m+2} = L_m - \frac{L_{m+1}}{m+1}.
\end{equation}

This derivation is valid for $m \geq 1$. We verify it holds for $m=0$: $L_2 = L_0 - L_1/1$. With $L_0=1$, $L_1 = \int_0^1 w(t)dt = \frac{\pi}{2}-1$, we get $L_2 = 1 - (\frac{\pi}{2}-1) = 2-\frac{\pi}{2}$. Direct calculation confirms this value.

For $g(m) = \sqrt{\frac{2}{\pi}}L_m$:
\begin{equation}\label{eq:g-discrete}
g(m+2) = g(m) - \frac{g(m+1)}{m+1}.
\end{equation}

\subsection{Arithmetic structure: the $\Q$-linear forms}

We now establish the explicit $\Q$-linear decomposition $K_m = p_m - \pi q_m$ with $p_m, q_m \in \Q_{\ge 0}$.

Define the sign-normalized quantity:
$$K_m = (-1)^m L_m.$$

From the recurrence $L_{m+2} = L_m - \frac{L_{m+1}}{m+1}$, multiplying by $(-1)^{m+2}$:
$$(-1)^{m+2}L_{m+2} = (-1)^{m+2}L_m - (-1)^{m+2}\frac{L_{m+1}}{m+1}.$$

This gives:
$$K_{m+2} = (-1)^{m+2}L_{m+2} = (-1)^m L_m + (-1)^{m+1} \frac{L_{m+1}}{m+1} = K_m + \frac{K_{m+1}}{m+1}.$$

Writing $K_m = p_m - \pi q_m$ with $p_m, q_m \in \Q_{\ge 0}$, the recurrence separates into:
\begin{align}
p_{m+2} &= p_m + \frac{p_{m+1}}{m+1}, \label{eq:p-rec}\\
q_{m+2} &= q_m + \frac{q_{m+1}}{m+1}. \label{eq:q-rec}
\end{align}

\textit{Initial conditions.} From $K_0 = L_0 = 1$: $p_0 = 1, q_0 = 0$. From $L_1 = \frac{\pi}{2} - 1$, we have $K_1 = -L_1 = 1 - \frac{\pi}{2}$: $p_1 = 1, q_1 = \frac{1}{2}$.

\textit{Alternation patterns.} The recurrences \eqref{eq:p-rec}--\eqref{eq:q-rec} with positive signs exhibit alternation patterns:

\begin{proposition}[Alternation and closed formulas]\label{prop:pq-formulas}
We establish explicit formulas for the rational coefficients $p_m$ and $q_m$ appearing in the $\Q$-linear decomposition $K_m = p_m - \pi q_m$. These coefficients exhibit alternation patterns and admit closed formulas in terms of double factorials.

For all $k \geq 1$:
\begin{enumerate}[label=(\roman*)]
\item $p_{2k+1} = p_{2k}$ (odd $p$-values freeze)
\item $q_{2k} = q_{2k-1}$ (even $q$-values freeze)
\item For even indices:
$$p_{2k} = \frac{(2k)\mskip-3mu\mathop{!!}}{(2k-1)\mskip-3mu\mathop{!!}}, \qquad q_{2k} = \frac{1}{2}\cdot\frac{(2k-1)\mskip-3mu\mathop{!!}}{(2k-2)\mskip-3mu\mathop{!!}}$$
\item For odd indices:
$$p_{2k+1} = \frac{(2k)\mskip-3mu\mathop{!!}}{(2k-1)\mskip-3mu\mathop{!!}}, \qquad q_{2k+1} = \frac{1}{2}\cdot\frac{(2k+1)\mskip-3mu\mathop{!!}}{(2k)\mskip-3mu\mathop{!!}}$$
\end{enumerate}
\end{proposition}

\begin{proof}
We establish the alternation patterns by induction, then derive the closed formulas via telescoping products. The key observation is that the positive sign in \eqref{eq:p-rec}--\eqref{eq:q-rec} creates a telescoping structure.

We first prove that $p_{2k+1} = p_{2k}$ for all $k \geq 1$.

Base case: Direct computation gives $p_0 = 1$, $p_1 = 1$, $p_2 = 1 + \frac{1}{1} = 2$, $p_3 = 1 + \frac{2}{2} = 2$. Thus $p_3 = p_2 = 2$.

Inductive step: Assume $p_{2k+1} = p_{2k}$. Then:
$$p_{2k+2} = p_{2k} + \frac{p_{2k+1}}{2k+1} = p_{2k}\left(1 + \frac{1}{2k+1}\right) = p_{2k} \cdot \frac{2k+2}{2k+1}.$$

And:
$$p_{2k+3} = p_{2k+1} + \frac{p_{2k+2}}{2k+2} = p_{2k} + \frac{1}{2k+2} \cdot p_{2k} \cdot \frac{2k+2}{2k+1} = p_{2k}\left(1 + \frac{1}{2k+1}\right) = p_{2k+2}.$$

Thus $p_{2k+3} = p_{2k+2}$, confirming the pattern.

Similarly, $q_{2k} = q_{2k-1}$ follows by identical induction, starting from $q_0 = 0$, $q_1 = \frac{1}{2}$, $q_2 = 0 + \frac{1/2}{1} = \frac{1}{2}$. So $q_2 = q_1 = \frac{1}{2}$.

We now derive closed formulas via telescoping products.
From the relation $p_{2k+2} = p_{2k} \cdot \frac{2k+2}{2k+1}$ derived above, we obtain a telescoping product for $k\ge 1$. Using $p_2 = 2$:
\[
p_{2k} = p_2 \prod_{j=1}^{k-1} \frac{2j+2}{2j+1} = 2 \cdot \frac{4 \cdot 6 \cdot 8 \cdots (2k)}{3 \cdot 5 \cdot 7 \cdots (2k-1)}.
\]
We express the terms using double factorials. The denominator is $(2k-1)\mskip-3mu\mathop{!!}$. The numerator is the product of even numbers from 4 to $2k$. Since $(2k)\mskip-3mu\mathop{!!} = 2 \cdot 4 \cdot 6 \cdots (2k)$, the numerator is $(2k)\mskip-3mu\mathop{!!}/2$. Thus:
\[
p_{2k} = 2 \cdot \frac{(2k)\mskip-3mu\mathop{!!}/2}{(2k-1)\mskip-3mu\mathop{!!}} = \frac{(2k)\mskip-3mu\mathop{!!}}{(2k-1)\mskip-3mu\mathop{!!}}.
\]
This establishes the formula for $p_{2k}$ in (iii). By the freezing property (i), $p_{2k+1}=p_{2k}$, establishing the formula in (iv).

For the $q$ coefficients, we use the relation $q_{2k+1} = q_{2k-1}\cdot\frac{2k+1}{2k}$ (derived similarly to $p$) with $q_1 = \frac{1}{2}$. The telescoping product is:
\[
q_{2k+1} = q_1 \prod_{j=1}^{k} \frac{2j+1}{2j} = \frac{1}{2} \cdot \frac{3 \cdot 5 \cdot 7 \cdots (2k+1)}{2 \cdot 4 \cdot 6 \cdots (2k)}.
\]
The numerator is $(2k+1)\mskip-3mu\mathop{!!}$ and the denominator is $(2k)\mskip-3mu\mathop{!!}$. Thus:
\[
q_{2k+1} = \frac{1}{2} \cdot \frac{(2k+1)\mskip-3mu\mathop{!!}}{(2k)\mskip-3mu\mathop{!!}}.
\]
This establishes the formula for $q_{2k+1}$ in (iv). By the freezing property (ii), $q_{2k}=q_{2k-1}$ (for $k\ge 1$). Applying the formula for odd indices to $q_{2k-1}$ yields the formula for $q_{2k}$ in (iii).
\end{proof}

\subsection{Analytic continuation and functional equation}

For $\Re z > 0$, extend $L_m$ (for $m \geq 1$) to 
$$L(z) = \int_0^1 zt^{z-1}\sqrt{\frac{1-t}{1+t}}\, \dd t.$$ 

The connection constant is $g(z) = \sqrt{\frac{2}{\pi}}L(z)$, which is holomorphic for $\Re z > 0$ (taking the principal branch of $t^{z-1}$).

\begin{lemma}[Growth bounds for $g$]\label{lem:growth-g}
The function $g(z)$ satisfies $|g(z)| = O(|z|^{-1/2})$ in vertical strips, ensuring the moderate growth required for Carlson's theorem. The proof, based on the explicit formula \eqref{eq:g-explicit} established below, is given in Section~\ref{sec:gamma-prototype}.
\end{lemma}
From the discrete recurrence \eqref{eq:g-discrete}, define the function 
$$h(z) = g(z+2) - g(z) + \frac{g(z+1)}{z+1}.$$ 

By \eqref{eq:g-discrete}, we have $h(m)=0$ for all $m\in\N$. By Lemma \ref{lem:growth-g}, $g(z)$ is holomorphic for $\Re z \geq \epsilon > 0$ and exhibits polynomial growth in vertical strips, hence $h(z)$ also has polynomial growth. This polynomial growth is of exponential type $<\pi$. Together with $h(n)=0$ for all $n\in\N$, Carlson's theorem (Theorem \ref{thm:carlson}) applies, yielding $h(z) \equiv 0$ and thus:
\begin{equation}\label{eq:g-AFE}
g(z+2) = g(z) - \frac{g(z+1)}{z+1}.
\end{equation}

\subsection{Explicit formula and verification}

We establish an explicit formula for $g(z)$ in terms of the Gamma function and prove its validity by verifying the functional equation.

\begin{proposition}\label{prop:g-explicit}
The connection constant $g(z)$ has simple poles at the negative integers $z \in \{-1,-2,-3,\ldots\}$ and is holomorphic at $z=0$. It is given by:
\begin{equation}\label{eq:g-explicit}
g(z) = \sqrt{2}\left[\frac{\Gamma\left(\frac{z}{2}+1\right)}{\Gamma\left(\frac{z+1}{2}\right)} - \frac{\Gamma\left(\frac{z+1}{2}\right)}{\Gamma\left(\frac{z}{2}\right)}\right].
\end{equation}
\end{proposition}

\begin{proof}
We prove this identity by showing that the proposed expression satisfies the functional equation \eqref{eq:g-AFE} and the initial conditions. Since both the integral definition and the proposed formula satisfy the same order-2 AFE with exponential type $<\pi$, agreement at $z=0$ and $z=1$ (verified below) propagates to all of $\Z$ via the recurrence. By Carlson's theorem, the difference of two functions with exponential type $<\pi$ that vanishes on $\N$ must vanish identically on $\C$.

We first verify that the proposed formula \eqref{eq:g-explicit} has exponential type $<\pi$ in vertical strips, as required for uniqueness. The uniform version of Stirling's formula shows that for Gamma ratios of the form $\Gamma(z/2+\alpha)/\Gamma(z/2+\beta)$, the dominant exponential factors $\exp(-\pi|\Im z|/4)$ cancel exactly, leaving polynomial growth (specifically, decay $O(|z|^{-1/2})$) in $|z|$. Since polynomial growth implies exponential type $0 < \pi$, the growth condition required for uniqueness is satisfied. (A detailed derivation is given in Section~\ref{sec:gamma-prototype}.)

Define the auxiliary function:
$$A(z) = \frac{\Gamma\left(\frac{z}{2}+1\right)}{\Gamma\left(\frac{z+1}{2}\right)}.$$

We have $A(z-1) = \frac{\Gamma\left(\frac{z-1}{2}+1\right)}{\Gamma\left(\frac{z-1+1}{2}\right)} = \frac{\Gamma\left(\frac{z+1}{2}\right)}{\Gamma\left(\frac{z}{2}\right)}$.
Thus, the proposed formula is $g(z) = \sqrt{2}(A(z) - A(z-1))$.

We derive relations between shifted values of $A(z)$ using $\Gamma(x+1)=x\Gamma(x)$:
\begin{align*}
A(z)A(z-1) &= \frac{\Gamma\left(\frac{z}{2}+1\right)}{\Gamma\left(\frac{z+1}{2}\right)} \cdot \frac{\Gamma\left(\frac{z+1}{2}\right)}{\Gamma\left(\frac{z}{2}\right)} = \frac{\frac{z}{2}\Gamma\left(\frac{z}{2}\right)}{\Gamma\left(\frac{z}{2}\right)} = \frac{z}{2}.
\end{align*}

This implies $A(z+1)A(z) = \frac{z+1}{2}$, so $A(z+1) = \frac{z+1}{2A(z)}$.
Also, $A(z+2) = \frac{z+2}{2A(z+1)} = \frac{z+2}{2} \frac{2A(z)}{z+1} = \frac{z+2}{z+1}A(z)$.

We now verify the functional equation $g(z+2) = g(z) - \frac{g(z+1)}{z+1}$. Dividing by $\sqrt{2}$, this becomes:
$$A(z+2)-A(z+1) = A(z)-A(z-1) - \frac{A(z+1)-A(z)}{z+1}.$$

Rearranging, we need to verify:
$$A(z+2)+A(z-1) = A(z+1)\frac{z}{z+1} + A(z)\frac{z+2}{z+1}.$$

We compute both sides using the relations derived above. Left-hand side:
$$A(z+2)+A(z-1) = \frac{z+2}{z+1}A(z) + \frac{z}{2A(z)}.$$

Right-hand side:
$$A(z+1)\frac{z}{z+1} + A(z)\frac{z+2}{z+1} = \frac{z+1}{2A(z)}\frac{z}{z+1} + \frac{z+2}{z+1}A(z) = \frac{z}{2A(z)} + \frac{z+2}{z+1}A(z).$$

Since the left-hand side equals the right-hand side, the functional equation is satisfied. Finally, we verify the initial values match those derived from the integral $L_m$. We have $L_1 = \pi/2-1$, hence 
$$g(1) = \sqrt{2/\pi}(\pi/2-1) = \frac{\pi-2}{\sqrt{2\pi}}.$$

We check the formula \eqref{eq:g-explicit} at $z=1$. 
$$A(1) = \frac{\Gamma(3/2)}{\Gamma(1)} = \sqrt{\pi}/2$$ 
and 
$$A(0) = \frac{\Gamma(1)}{\Gamma(1/2)} = 1/\sqrt{\pi}.$$

Then 
$$ \sqrt{2}(A(1)-A(0)) = \sqrt{2}(\sqrt{\pi}/2 - 1/\sqrt{\pi}) = \sqrt{2}\frac{\pi-2}{2\sqrt{\pi}} = \frac{\pi-2}{\sqrt{2\pi}}=g(1).$$
This matches.  We also verify a second initial value to ensure complete uniqueness. For $z=0$: $L_0 = 1$, hence 
$$g(0) = \sqrt{2/\pi} \cdot 1 = \sqrt{2/\pi}.$$

Using the explicit formula \eqref{eq:g-explicit}, we have
$$g(0) = \sqrt{2}(A(0)-A(-1)).$$
We compute
$$A(0) = \frac{\Gamma(1)}{\Gamma(1/2)} = \frac{1}{\sqrt{\pi}}$$
and
$$A(-1) = \lim_{z \to -1} A(z) = \lim_{z \to -1} \frac{\Gamma(z/2+1)}{\Gamma((z+1)/2)} = 0.$$
Thus
$$g(0) = \sqrt{2}(1/\sqrt{\pi} - 0) = \sqrt{2/\pi},$$
confirming the expected value. Since the expression \eqref{eq:g-explicit} satisfies the same linear functional equation as the connection constant $g(z)$, matches the initial conditions, and exhibits the required growth properties, the identity is proven by uniqueness.
\end{proof}

This formula expresses $g(z)$ as a combination of ratios of ordinary Gamma functions, confirming its regular nature and completing the proof of Theorem \ref{thm:pi-world}.

\section{The AGF prototypes: $\Gamma$, $f$, and $g$}\label{sec:gamma-prototype}

Having established the explicit formulas for $f(z)$ and $g(z)$ in Sections~\ref{sec:e-world} and~\ref{sec:pi-world}, we now reveal the unifying structure: these functions, together with Euler's Gamma function, are instances of a single class characterized by three properties. We verify these properties systematically, completing the proofs of the growth bounds invoked earlier.

\subsection{The three defining properties}

Recall from Definition~\ref{def:AGF-intro} that an additive Gamma function satisfies:
\begin{enumerate}[label=(\roman*)]
\item An additive functional equation with rational coefficients;
\item Moderate growth in vertical strips (polynomial, hence exponential type $<\pi$);
\item Uniqueness via normalization on an interval of length equal to the order.
\end{enumerate}

We now verify these properties for $\Gamma(z)$ (order 1), $f(z)$ (order 2), and $g(z)$ (order 2).

\subsection{Euler's Gamma function: the order-1 prototype}

Euler's product formula (1729), refined by Gauss, can be written:
$$\Gamma(z) = \lim_{n\to\infty} \frac{n! n^z}{z(z+1)\cdots(z+n)}.$$

This classical formula is typically presented as a limit defining $\Gamma(z)$. We now show its deeper structure: \textit{it is the connection constant of a first-order P-recursive sequence}.

Consider the sequence
\begin{equation}\label{eq:gamma-recurrence}
w_{n+1}(z) = \frac{n+1}{n+z} w_n(z), \quad w_1(z) = 1.
\end{equation}

The solution is $w_n(z) = \frac{n!}{(z)_n}$ where $(z)_n = z(z+1)\cdots(z+n-1)$ is the Pochhammer symbol. Using Stirling's approximation, $w_n(z) \sim \Gamma(z) n^{1-z}$, identifying $\Gamma(z)$ as the connection constant.

\textbf{Verification of AGF properties for $\Gamma(z)$:}
\begin{enumerate}[label=(\roman*)]
\item \textbf{AFE:} $\Gamma(z+1) = z\Gamma(z)$, with rational coefficient $R_0(z) = z$, $R_1(z) = -1$.
\item \textbf{Growth:} By Stirling's formula, $|\Gamma(z)| \sim \sqrt{2\pi}|z|^{\Re z - 1/2}e^{-\pi|\Im z|/2}$ in vertical strips. The exponential factor has exponent $\pi/2 < \pi$.
\item \textbf{Uniqueness:} $\Gamma(z)$ is determined by the AFE together with $\Gamma(1) = 1$.
\end{enumerate}

The holonomic triangle for $\Gamma(z)$ is illustrated in Figure~\ref{fig:gamma-triangle}.

\subsection{Growth bounds for $f(z)$: proof of Lemma~\ref{lem:growth-f}}

We use the explicit formula $f(z) = \ee^{-1-i\pi z}\gamma(z+2,-1)$ from \eqref{eq:f-explicit}. For the incomplete Gamma function $\gamma(a,x)$ with fixed $x$ and large $|a|$, according to \cite[Eq.~8.11.6]{DLMF}:
\[
\gamma(a,x) = \frac{x^a \ee^{-x}}{a} \left(1 + O\left(\frac{1}{a}\right)\right)
\]
uniformly in the sector $|\arg a| \le \pi-\delta$.

Setting $a=z+2$ and $x=-1$, using the principal branch $(-1)^{z+2} = \ee^{i\pi z}$ and $\ee^{-(-1)} = \ee$:
\[
\gamma(z+2,-1) = \frac{\ee^{i\pi z+1}}{z+2} \left(1 + O\left(\frac{1}{z}\right)\right).
\]
Substituting:
\[
f(z) = \ee^{-1-i\pi z} \cdot \frac{\ee^{i\pi z+1}}{z+2} \left(1 + O\left(\frac{1}{z}\right)\right) = \frac{1}{z+2} + O\left(\frac{1}{z^2}\right).
\]
Therefore $f(z) = O(1/|z|)$ exhibits polynomial decay in vertical strips, satisfying the growth bound with exponential type $A=0 < \pi$. \qed

\textbf{Verification of AGF properties for $f(z)$:}
\begin{enumerate}[label=(\roman*)]
\item \textbf{AFE:} $f(z+2) = (z+2)f(z) - (z+2)f(z+1)$, with rational coefficients.
\item \textbf{Growth:} Polynomial decay $O(1/|z|)$ as shown above.
\item \textbf{Uniqueness:} Determined by the AFE together with $f(0) = 1/\ee$ and $f(1)$.
\end{enumerate}

\subsection{Growth bounds for $g(z)$: proof of Lemma~\ref{lem:growth-g}}

We use the explicit formula $g(z) = \sqrt{2}(A(z) - A(z-1))$ from \eqref{eq:g-explicit}, where $A(z) = \Gamma(\frac{z}{2}+1)/\Gamma(\frac{z+1}{2})$. For Gamma ratios with large argument, according to \cite[Eq.~5.11.12]{DLMF}:
\[
\frac{\Gamma(w+a)}{\Gamma(w+b)} = w^{a-b} \left[ 1 + \frac{(a-b)(a+b-1)}{2w} + O\left(\frac{1}{w^2}\right) \right].
\]

Setting $w=z/2$:
\[
A(z) = \frac{\Gamma(w+1)}{\Gamma(w+1/2)} = w^{1/2} \left[ 1 + \frac{1}{8w} + O\left(\frac{1}{w^2}\right) \right],
\]
\[
A(z-1) = \frac{\Gamma(w+1/2)}{\Gamma(w)} = w^{1/2} \left[ 1 - \frac{1}{8w} + O\left(\frac{1}{w^2}\right) \right].
\]
Therefore:
\[
g(z) = \sqrt{2}(A(z)-A(z-1)) = \sqrt{2} w^{1/2} \cdot \frac{1}{4w} + O(w^{-3/2}) = O(|z|^{-1/2}).
\]
Thus $g(z)$ exhibits polynomial decay in vertical strips, satisfying the growth bound with exponential type $A=0 < \pi$. \qed

\textbf{Verification of AGF properties for $g(z)$:}
\begin{enumerate}[label=(\roman*)]
\item \textbf{AFE:} $g(z+2) = g(z) - \frac{g(z+1)}{z+1}$, with rational coefficients.
\item \textbf{Growth:} Polynomial decay $O(|z|^{-1/2})$ as shown above.
\item \textbf{Uniqueness:} Determined by the AFE together with $g(0) = \sqrt{2/\pi}$ and $g(1)$.
\end{enumerate}

\subsection{Summary: the AGF hierarchy}

We have verified that $\Gamma(z)$, $f(z)$, and $g(z)$ satisfy all three defining properties of additive Gamma functions:

\begin{center}
\begin{tabular}{lccc}
\hline
Function & Order & AFE & Growth \\
\hline
$\Gamma(z)$ & 1 & $\Gamma(z+1) = z\Gamma(z)$ & $O(|z|^{\Re z - 1/2}e^{-\pi|\Im z|/2})$ \\
$f(z)$ & 2 & $f(z+2) = (z+2)[f(z) - f(z+1)]$ & $O(|z|^{-1})$ \\
$g(z)$ & 2 & $g(z+2) = g(z) - \frac{g(z+1)}{z+1}$ & $O(|z|^{-1/2})$ \\
\hline
\end{tabular}
\end{center}

This confirms that the AGF framework captures a genuine class of special functions, with Euler's $\Gamma(z)$ as the founding example and $f(z)$, $g(z)$ as order-2 instances exhibiting the irregular/regular dichotomy.

\section{The structure of additive Gamma functions}\label{sec:general-theory}

\subsection{From examples to general theory}

The functions $\Gamma(z)$, $f(z)$, and $g(z)$ share three properties: connection constant origin, additive functional equations with rational coefficients, and controlled growth. This motivates a general definition.

We call these functions \textbf{additive Gamma functions} because, like Euler's Gamma function (the order-1 prototype), they are uniquely determined by a linear functional equation together with a growth condition, and emerge from asymptotic analysis of holonomic sequences.

The key mechanism ensuring rational coefficients is the integer slope condition. Consider a Gamma factor $\Gamma(n+\alpha z+\beta)$ in $u_n(z) \sim \Lambda(n,z)h(z)$. Under $z\mapsto z+1$:
$$\frac{\Gamma(n+\alpha(z+1)+\beta)}{\Gamma(n+\alpha z+\beta)} = \frac{\Gamma(n+\alpha z + \alpha + \beta)}{\Gamma(n+\alpha z+\beta)}.$$

When $\alpha \in \Z$ with $\alpha > 0$, this ratio equals $(n+\alpha z+\beta)_\alpha$, polynomial in $z$. When $\alpha \in \Z$ with $\alpha < 0$, it equals $\frac{1}{(n+\alpha z+\beta-\alpha)_{|\alpha|}}$, rational in $z$. However, if $\alpha\notin\Z$, the poles form two disjoint infinite arithmetic progressions, preventing rationality.

\subsection{The Integer Slope Condition}

We now state and prove the central technical tool:

\begin{lemma}[Integer Slope Condition]\label{lem:integer-slope}
Let $\phi(z) = \alpha z + \beta$ be a linear function with $\alpha, \beta \in \C$ and $\alpha \neq 0$. The ratio
$$R(z) := \frac{\Gamma(\phi(z+1))}{\Gamma(\phi(z))} = \frac{\Gamma(\alpha z + \alpha + \beta)}{\Gamma(\alpha z + \beta)}$$
is a rational function of $z$ if and only if $\alpha \in \Z$.
\end{lemma}

\begin{proof}
We treat sufficiency and necessity separately.

\textit{Sufficiency ($\Leftarrow$).} Assume $\alpha \in \Z$. We distinguish two cases.

If $\alpha = k > 0$, using the functional equation $\Gamma(w+1) = w\Gamma(w)$ iteratively $k$ times yields $\Gamma(w+k) = (w)_k \Gamma(w)$, where $(w)_k = w(w+1)\cdots(w+k-1)$ is the Pochhammer symbol. Setting $w = \alpha z + \beta$, we obtain
$$R(z) = \frac{\Gamma(w+k)}{\Gamma(w)} = (w)_k = \prod_{j=0}^{k-1}(\alpha z + \beta + j),$$
a polynomial in $z$ of degree $k$.

If $\alpha = -k < 0$ with $k > 0$, let $w = \alpha z + \beta$. The numerator is $\Gamma(w-k)$. Writing $\Gamma(w) = \Gamma((w-k)+k)$ and applying the product formula, we get $\Gamma(w) = (w-k)_k \Gamma(w-k)$. Thus
$$R(z) = \frac{\Gamma(w-k)}{\Gamma(w)} = \frac{1}{(w-k)_k} = \frac{1}{\prod_{j=0}^{k-1}(\alpha z + \beta - k + j)},$$
the reciprocal of a polynomial, hence rational.

\textit{Necessity ($\Rightarrow$).} Assume $R(z) \in \Q(z)$. We prove $\alpha \in \Z$ by contradiction.

A rational function has finitely many poles and zeros. The function $\Gamma$ has no zeros and simple poles at $\{0, -1, -2, \ldots\}$.

The poles of the denominator $\Gamma(\alpha z + \beta)$ occur at $\alpha z + \beta = -k$ for $k \in \N$, giving
$$P_{\mathrm{denom}} = \left\{ z_k = \frac{-\beta - k}{\alpha} : k \in \N \right\}.$$

The poles of the numerator $\Gamma(\alpha z + \alpha + \beta)$ occur at
$$P_{\mathrm{numer}} = \left\{ z'_k = \frac{-\beta - \alpha - k}{\alpha} = z_k - 1 : k \in \N \right\}.$$

Since $\Gamma$ has no zeros, the numerator is finite and nonzero at every point of $P_{\mathrm{denom}}$ that does not belong to $P_{\mathrm{numer}}$. At such points, $R(z)$ has a zero (from the pole of the denominator).

If $\alpha \notin \Z$, the sets $P_{\mathrm{denom}}$ and $P_{\mathrm{numer}}$ are disjoint: equality $z_j = z'_k$ would imply $\alpha = k - j \in \Z$, contradicting our assumption. Consequently, $R(z)$ vanishes at every point of the infinite set $P_{\mathrm{denom}}$.

The only rational function with infinitely many zeros is identically zero. But $R(z) \not\equiv 0$ since $\Gamma$ is nowhere zero. This contradiction establishes $\alpha \in \Z$.
\end{proof}

\subsection{Characterization of additive Gamma functions}

As motivated in the introduction (Definition \ref{def:AGF-intro}), the properties common to $\Gamma(z)$, $f(z)$, and $g(z)$ lead to the formal definition of an additive Gamma function: a meromorphic function satisfying a homogeneous additive functional equation with rational coefficients and exhibiting polynomial growth in vertical strips. This definition captures the essential structure shared by all connection constants arising from holonomic sequences with integer-slope asymptotics.

\begin{theorem}[Characterization of AGFs]\label{thm:AGF-equivalence}
A meromorphic function $h$ is an AGF of order $r$ if and only if $h$ arises as a connection constant from a sequence $u_n(z)$ that is holonomic in $(n,z)$ with integer-slope asymptotics:
$$u_n(z) \sim \Lambda(n,z)h(z) \quad \text{as } n \to \infty,$$
where
$$\Lambda(n,z) = \lambda^n n^{\rho(z)} \prod_j \Gamma(n+\alpha_j z+\beta_j)^{m_j}$$
with constant $\lambda\in\C^*$, affine $\rho(z)$ with integer slope, integer slopes $\alpha_j\in\Z$, and $m_j\in\Z$.
\end{theorem}

\begin{proof}
\textit{Direction ($\Leftarrow$).} Suppose $h(z)$ is the connection constant of a holonomic sequence $u_n(z)$ with integer-slope asymptotics $u_n(z) \sim \Lambda(n,z)h(z)$. By Proposition~\ref{prop:prec-to-afe}, the holonomicity guarantees the existence of a contiguity relation in $z$, from which $h(z)$ inherits an additive functional equation with rational coefficients. The moderate growth condition follows from the structure of $\Lambda(n,z)$: Gamma factors with integer slopes have ratios (under parameter shifts in $z$) that exhibit polynomial growth via Stirling estimates, as demonstrated for $f(z)$ and $g(z)$ in Section~\ref{sec:gamma-prototype}. Hence $h$ is an AGF.

\textit{Direction ($\Rightarrow$).} Suppose $h$ is an AGF of order $r$. We construct a sequence $u_n(z)$ holonomic in $(n,z)$ having $h(z)$ as its connection constant.

A direct definition $u_n(z) := n^z h(z)$ is not valid because $n^z$ is not P-recursive in $n$ over $\Q(n,z)$: the ratio $(n+1)^z/n^z = (1+1/n)^z$ is not rational in $n$.

Instead, we construct a holonomic sequence asymptotically equivalent to $n^z$. The sequence
$$w_n(z) := \frac{n!}{\Gamma(n+1-z)}$$
satisfies recurrences in both variables:
$$w_{n+1}(z) = \frac{n+1}{n+1-z} w_n(z), \qquad w_n(z+1) = (n-z) w_n(z),$$
with coefficients in $\Q(n,z)$, so $w_n(z)$ is holonomic in $(n,z)$. By Stirling's formula, $w_n(z) \sim n^z$ as $n \to \infty$.

More generally, the class of holonomic sequences is closed under multiplication \cite{Stanley1999}. The components of a general shell $\Lambda(n,z)$ are:
\begin{enumerate}
\item The geometric term $\lambda^n$ satisfies $a_{n+1} = \lambda \cdot a_n$ (holonomic, trivial in $z$).
\item The power term $n^{\rho(z)}$ with $\rho(z) = c_1 z + c_0$ is asymptotically equivalent to $w_n(z)^{c_1} \cdot n^{c_0}$, a product of holonomic sequences.
\item Each Gamma factor $\Gamma(n + \alpha_j z + \beta_j)^{m_j}$ with $\alpha_j \in \Z$ is holonomic in $(n,z)$ via the functional equations of $\Gamma$.
\end{enumerate}

Let $\widetilde{\Lambda}(n,z)$ be the product of these holonomic components. Then $\widetilde{\Lambda}(n,z)$ is holonomic in $(n,z)$ and $\widetilde{\Lambda}(n,z) \sim \Lambda(n,z)$ as $n \to \infty$.

Define $u_n(z) := \widetilde{\Lambda}(n,z) \cdot h(z)$. Since $h(z)$ is independent of $n$, the sequence $u_n(z)$ inherits holonomicity from $\widetilde{\Lambda}(n,z)$ and has the required asymptotics $u_n(z) \sim h(z) \Lambda(n,z)$.
\end{proof}

\begin{proposition}[From Holonomic Sequence to AFE]\label{prop:prec-to-afe}
Let $u_n(z)$ be holonomic in $(n,z)$, i.e., satisfy compatible linear recurrences in both $n$ and $z$ with coefficients in $\Q(n,z)$. Assume integer-slope asymptotics as in Theorem~\ref{thm:AGF-equivalence}:
\begin{equation*}
u_n(z) \sim \Lambda(n,z)h(z) \quad \text{where} \quad 
\Lambda(n,z) = \lambda^n n^{\rho(z)} \prod_j \Gamma(n+\alpha_j z+\beta_j)^{m_j}
\end{equation*}
with $\rho(z)$ affine of integer slope and $\alpha_j \in \Z$.

Then the connection constant $h(z)$ satisfies an AFE with rational 
coefficients.
\end{proposition}

\begin{proof}
Since $u_n(z)$ is holonomic in $(n,z)$, elimination theory in Ore algebras \cite{ChyzakSalvy1998} guarantees the existence of a contiguity relation in $z$:
\begin{equation}\label{eq:contiguity}
\sum_{k=0}^{r'} S_k(n,z) u_n(z+k) = 0,
\end{equation}
where $S_k(n,z) \in \Q(n,z)$ are rational functions.

Substituting the asymptotic form $u_n(z+k) \sim \Lambda(n,z+k)h(z+k)$ and dividing by $\Lambda(n,z)$ yields
\begin{equation}\label{eq:asymp-contiguity}
\sum_{k=0}^{r'} h(z+k) \cdot T_k(n,z) \approx 0, \quad \text{where } T_k(n,z) := S_k(n,z) \frac{\Lambda(n,z+k)}{\Lambda(n,z)}.
\end{equation}

We analyze the ratio $\Lambda(n,z+k)/\Lambda(n,z)$. The power part contributes $n^{c_1 k}$ since $\rho(z) = c_1 z + c_0$ is affine with $c_1 \in \Z$; this is a polynomial in $n$. For each Gamma factor, the ratio
$$G_j(n,z,k) = \frac{\Gamma(n+\alpha_j(z+k)+\beta_j)}{\Gamma(n+\alpha_j z+\beta_j)}$$
equals $(n+\alpha_j z+\beta_j)_M$ if $M = \alpha_j k > 0$, or $1/(n+\alpha_j z+\beta_j+M)_{|M|}$ if $M < 0$. By Lemma~\ref{lem:integer-slope}, since $\alpha_j \in \Z$, these are rational in $(n,z)$.

Thus each $T_k(n,z)$ admits an expansion in powers of $n$:
$$T_k(n,z) = n^{\sigma_k} \left( R_k^{(0)}(z) + R_k^{(1)}(z) n^{-1} + O(n^{-2}) \right)$$
where $\sigma_k \in \Z$ and $R_k^{(j)}(z) \in \Q(z)$.

Let $\sigma = \max_k \sigma_k$ be the dominant exponent. Dividing \eqref{eq:asymp-contiguity} by $n^\sigma$ and taking $n \to \infty$, the leading coefficient must vanish:
$$ \sum_{k : \sigma_k = \sigma} h(z+k) \cdot R_k^{(0)}(z) = 0. $$

If this equation is trivial (all $R_k^{(0)} = 0$ for the maximal indices), we proceed to the next power of $n$. Since \eqref{eq:contiguity} holds exactly for all $n$, some coefficient must be nontrivial, yielding a nondegenerate AFE
$$\sum_{k=0}^{r'} R_k(z) h(z+k) = 0, \qquad R_k(z) \in \Q(z). \qedhere$$
\end{proof}

\begin{example}[Derivation of AFE for $f(z)$]\label{ex:afe-derivation}
We illustrate Proposition~\ref{prop:prec-to-afe} with the $\ee$-world example.

The sequence $u_n(z)$ satisfies the P-recurrence (after clearing denominators):
\begin{equation}\label{eq:prec-example}
(n+z)u_{n+2}(z) - (n+z)u_{n+1}(z) - u_n(z) = 0.
\end{equation}

The functional equation for $f(z)$ does not arise directly from this recurrence, which relates terms at fixed $z$. Instead, one derives a contiguity relation linking $u_n(z)$, $u_n(z+1)$, and $u_n(z+2)$ via elimination. Shifting $z \mapsto z+1$ in \eqref{eq:prec-example} yields
$$
(n+z+1)u_{n+2}(z+1) - (n+z+1)u_{n+1}(z+1) - u_n(z+1) = 0,
$$
and similarly for $z \mapsto z+2$. The three equations (at $z$, $z+1$, $z+2$) form a system from which a contiguity relation in $z$ can be extracted.

With the asymptotic ansatz $u_n(z) \sim f(z) \cdot n$, we have $u_n(z+k) \sim f(z+k) \cdot n$. Since the asymptotic shell $\Lambda(n,z) = n$ is independent of $z$, the ratio $\Lambda(n,z+k)/\Lambda(n,z) = 1$ contributes no $z$-dependence.

The leading-order substitution into the contiguity relation yields a degenerate equation (the $O(n)$ terms cancel). To obtain the correct AFE, one uses the refined expansion $u_n(z) = f(z) \cdot n + c(z) + O(1/n)$ and extracts the coefficient of the leading non-vanishing term. Matching at order $O(1)$ recovers the functional equation
$$f(z+2) = (z+2)[f(z) - f(z+1)]$$
established directly in Section~\ref{sec:e-world}.

This example shows that while the general mechanism of Proposition~\ref{prop:prec-to-afe} applies, extracting the AFE may require asymptotic refinement beyond leading order when the shell $\Lambda(n,z)$ has trivial $z$-dependence.
\end{example}

\subsection{Examples}

We have three fundamental examples:
\begin{enumerate}
\item $\Gamma(z)$: order 1, $\Lambda(n,z)=n^{1-z}$, equation $\Gamma(z+1)=z\Gamma(z)$
\item $f(z)$: order 2, $\Lambda(n,z)=n$, equation \eqref{eq:f-AFE}
\item $g(z)$: order 2, $\Lambda(n,z)=\sqrt{n}$, equation \eqref{eq:g-AFE}
\end{enumerate}

The additive paradigm contrasts with multiplicative generalizations (Barnes G-function \cite{Barnes1901}). Our AGFs combine linearly at shifted arguments, generalizing the linear recurrence structure defining $\Gamma(z)$.

\section{The regular-irregular dichotomy}\label{sec:dichotomy}

\subsection{The fundamental distinction}

Formulas \eqref{eq:f-explicit} and \eqref{eq:g-explicit} reveal a structural difference rooted in asymptotic behavior at infinity.

For $g(z)$: $g(z+2) = g(z) - \frac{g(z+1)}{z+1}$. As $|z| \to \infty$, the coefficient $\frac{1}{z+1} \to 0$, making this \textit{regular at infinity} \cite{BirkhoffTrjitzinsky1933}. Regular equations have solutions through ordinary hypergeometric functions, reducing to Gamma ratios. Formula \eqref{eq:g-explicit} confirms this.

For $f(z)$: $f(z+2) = (z+2)f(z) - (z+2)f(z+1)$. Coefficients grow linearly, making this \textit{irregular at infinity}. Irregular equations require confluent hypergeometric functions. Formula \eqref{eq:f-explicit} involves $\gamma(z+2,-1)$, confirming this.

\subsection{Implications}

This dichotomy shows AGFs encompass both regular and irregular cases, richer than Gamma ratios alone. The mirror symmetry \eqref{eq:mirror} establishes correspondence between $\ee$ and $\pi$, and between ordinary and confluent hypergeometric realms.

\section{The holonomic perspective}\label{sec:holonomic}

\subsection{Scope and limitations}

An additive Gamma function of order $r$ is a single mathematical object that manifests in three equivalent forms (Figure~\ref{fig:triangle}): as a P-recursive sequence with rational coefficients and integer slopes, as an additive functional equation for the connection constant, and as a D-finite differential equation for the generating function. This correspondence applies specifically to sequences with factorisable asymptotics having integer slopes, ensuring rational coefficients in the functional equations.

\subsection{The holonomic triangle of order-$r$ equations}

\begin{principle}[The holonomic triangle]\label{prin:holonomic}
This principle provides a conceptual framework positioning AGFs within the broader landscape of holonomic objects. It is not a formal theorem but rather an organizing perspective that guides the theory.

For P-recursive sequences with integer-slope asymptotics and moderate algebraic growth, an additive Gamma function of order $r$ typically admits three related descriptions:

\begin{enumerate}
\item \textbf{P-recurrence:} A sequence $u_n(z)$ satisfying
$$\sum_{k=0}^r p_k(n,z) u_{n+k}(z) = 0$$
with $p_k(n,z) \in \Q(n,z)$ and integer-slope asymptotics $u_n(z) \sim \Lambda(n,z)h(z)$.

\item \textbf{AFE:} The connection constant $h(z)$ satisfying
$$\sum_{k=0}^r R_k(z) h(z+k) = 0$$
with rational coefficients $R_k(z) \in \Q(z)$.

\item \textbf{D-finite ODE:} The ordinary generating function $U(x,z) = \sum_n u_n(z)x^n$ satisfying a linear ODE:
$$\sum_{j=0}^m Q_j(x,z) \frac{\dd^j U}{\dd x^j} = 0,$$
with $Q_j(x,z) \in \Q(x,z)$. Under the integer-slope and moderate growth conditions, the order $m$ depends on both the width $r$ of the P-recurrence and the maximum degree $d$ of its polynomial coefficients (see Remark~\ref{rem:order-ode}). Based on elimination theory in Ore algebras \cite{ChyzakSalvy1998,Koutschan2010}, the order is typically at most $(d+1)r$, though often much lower in practice.
\end{enumerate}

The correspondences are:
\begin{itemize}
\item (1) $\to$ (2): Via Proposition \ref{prop:prec-to-afe}, by deriving the contiguity relation in $z$ for $u_n(z)$ and taking the asymptotic limit $n\to\infty$, utilizing the integer slope condition to ensure rationality.
\item (1) $\to$ (3): Standard generating function techniques \cite{FlajoletSedgewick2009}.
\item (2) $\leftrightarrow$ (3): Via singularity analysis and the transfer theorem.
\item (2) $\to$ (1): The connection constant $h(z)$ of an AGF determines a canonical P-recursive sequence through its generating function, as shown in Theorem~\ref{thm:AGF-equivalence}.
\end{itemize}

We emphasize that AGFs are not D-finite functions in general; rather, they are organized by additive functional equations. In favorable cases (such as our examples $f(z)$ and $g(z)$), they admit hypergeometric or Gamma-ratio representations that connect to the D-finite world.

This framework connects to Zeilberger's holonomic systems approach \cite{Zeilberger1990} and creative telescoping methods \cite{Zeilberger1991}. Modern computational tools based on Ore algebras \cite{ChyzakSalvy1998,Koutschan2010,KauersPaule2011} make these correspondences constructive and algorithmic.
\end{principle}

For P-recursive sequences satisfying the integer slope condition, the same mathematical object of order $r$ appears in three forms, constituting the holonomic triangle (Figure~\ref{fig:triangle}):

\begin{figure}[ht]
\centering
\begin{tikzcd}[row sep=2.5cm, column sep=2cm]
& \boxed{
\begin{tabular}{c}
\textbf{P-recurrence} \\[4pt]
order $r$
\end{tabular}}
\arrow[dl, <->, "\text{  $\,\,\,\,\,\,\,\,\,\,\,\,\,\,\,\,\,\,\,\,\,\,\,\,\,\,\,$       Coefficients mapping}" near end, sloped] 
\arrow[dr, <->, "\text{$\,\,\,\,\,\,\,\,\,\,\,\,\,\,\,\,\,\,\,\,\,\,\,\,\,\,\,\,\,$ Generating function}" near start, sloped] 
& \\
\boxed{
\begin{tabular}{c}
\textbf{AFE} \\[4pt]
order $r$
\end{tabular}}
\arrow[rr, <->, "\text{$\,\,\,\,\,\,\,\,\,\,\,\,\,\,\,\,\,\,\,\,\,\,\,\,\,\,\,\,\,\,\,\,\,\,\,\,\,\,\,\,\,\,\,\,\,\,\,\,\,\,\,\,\,\,\,$ Singularity analysis of the ODE solution}" near start] 
& & 
\boxed{
\begin{tabular}{c}
\textbf{D-finite ODE} \\[4pt]
order $\leq (d+1)r$
\end{tabular}}
\end{tikzcd}
\vspace{0.5em}
\caption{The holonomic triangle. The vertices represent the P-recurrence (order $r$), the AFE (order $r$), and the D-finite ODE (order $\leq (d+1)r$ for the ordinary generating function, where $d$ is the maximum degree of polynomial coefficients in the P-recurrence). Note: the bound $(d+1)r$ is worst-case; in practice the ODE order is often much lower---for instance, our order-2 examples yield ODEs of order 1.}
\label{fig:triangle}
\end{figure}

\subsection{Concrete holonomic certificates}

The holonomic triangle is not merely a conceptual framework—it can be made explicit through concrete differential equations. We illustrate this with our order-2 examples.

\begin{proposition}[Holonomic certificate for the $\ee$-world]\label{prop:GF-ODE-e}
For each fixed parameter $z\in\C$, the generating function $U_z(x) := \sum_{n\ge1}u_n(z)x^n$ satisfies a first-order linear differential equation with rational coefficients in $x$:
\[
  (1-x)U'_z(x) = U_z(x) \left(\frac{x^2-x+2-z+zx}{x}\right) + zx.
\]
This equation, derived in Section~\ref{sec:e-world}, makes the holonomic structure explicit and verifiable.
\end{proposition}

\begin{proposition}[Holonomic certificate for the $\pi$-world]\label{prop:GF-ODE-pi}
For each fixed parameter $z\in\C$, the generating function $V_z(x) := \sum_{n\ge1}v_n(z)x^n$ satisfies a first-order linear differential equation with rational coefficients in $x$:
\[
  (1-x^2)V_z'(x) + V_z(x)\left(\frac{z-2}{x}-zx-1\right) = zx.
\]
This equation, derived in Section~\ref{sec:pi-world}, provides the D-finite vertex of the holonomic triangle for $g(z)$.
\end{proposition}

These explicit ODEs make the connection between the three vertices of the holonomic triangle completely constructive. Given the P-recurrences for $u_n(m)$ and $v_n(m)$, standard generating function manipulations yield these ODEs. Conversely, singularity analysis of the ODE solutions recovers the coefficient asymptotics and the AFEs for $f(z)$ and $g(z)$.

\begin{remark}[Algorithmic verification]
Modern computer algebra systems (e.g., Maple with \texttt{gfun}, Mathematica, or SageMath) can verify these certificates algorithmically: starting from the P-recurrences, elimination techniques produce the ODEs; solving the ODEs and extracting singularity expansions confirms the asymptotic formulas; and the AFEs can be verified numerically or symbolically using the explicit formulas \eqref{eq:f-explicit} and \eqref{eq:g-explicit}.
\end{remark}

\begin{figure}[ht]
\centering
\begin{tikzcd}[row sep=2.5cm, column sep=2.3cm]
  & {\boxed{\begin{array}{c}
      \text{\bfseries P-recurrence}\\[4pt]
      w_{n+1}(z)=\dfrac{n+1}{n+z}\,w_n(z)\\[2pt]
      w_1(z)=1
    \end{array}}}
    \arrow[dl,<->,sloped,"{\,z\;\leftrightarrow\;\dfrac{n+1}{n+z}\,}"{pos=.5}]
    \arrow[dr,<->,sloped,"{\,W(x,z)=\sum_{n=1}^\infty w_n(z)\,x^n\,}"{pos=.5}]
  & \\
  {\boxed{\begin{array}{c}
      \text{\bfseries AFE}\\[4pt]
      \Gamma(z+1)=z\,\Gamma(z)\\[2pt]
      \Gamma(1)=1
    \end{array}}}
    \arrow[rr,<->,
      "{\substack{\text{Singularity analysis of }\\[2pt]
      W(x,z)=\,z\bigl({}_2F_1(1,1;z;x)-1\bigr)}}"{pos=.5}]
  & &
  {\boxed{\begin{array}{c}
      \text{\bfseries D\text{-}finite ODE}\\[4pt]
      (1-x)W' + \dfrac{z-1-x}{x}\,W = z
    \end{array}}}
\end{tikzcd}
\vspace{0.5em}
\caption{The holonomic triangle for Euler's Gamma function, illustrating the three-way equivalence between P-recurrence, AFE, and D-finite ODE under the integer slope condition.}
\label{fig:gamma-triangle}
\end{figure}

Euler's Gamma function provides the concrete order-1 illustration of this triangle (Figure~\ref{fig:gamma-triangle}). Our examples $f(z)$ and $g(z)$ have trivial integer slopes since $\Lambda(n,z) = n$ and $\Lambda(n,z) = \sqrt{n}$ involve no Gamma factors with $z$-dependence.

\section{Concluding remarks}

This work establishes mirror symmetry between recurrences connecting $\ee$ and $\pi$. Through parametrization and asymptotic analysis, connection constants extend to meromorphic functions satisfying additive functional equations with rational coefficients. The arithmetic duality \eqref{eq:duality} between $\Z$-linear forms in $\ee$ and $\Q$-linear forms in $\pi$ provides the motivating structure.

We constructed order-2 AGFs $f(z)$ and $g(z)$, recognizing Euler's Gamma function as the order-1 prototype. The Integer Slope Condition (Lemma~\ref{lem:integer-slope}) ensures rational coefficients, and Theorem~\ref{thm:AGF-equivalence} establishes the equivalence between AGFs and connection constants from integer-slope P-recurrences. The explicit formulas \eqref{eq:f-explicit} and \eqref{eq:g-explicit} reveal the regular/irregular dichotomy.

\subsection{Open questions}

The integer slope condition delineates where rational-coefficient AFEs emerge. For sequences with non-integer slopes or non-factorisable asymptotics, whether functional equations exist (and with what coefficients) remains open; beyond this boundary, transcendental coefficients would likely be required. Several directions merit further investigation:

\begin{enumerate}
\item \textbf{Classification at order 2.} Systematic study of order-2 AGFs and their connection to Heun's differential equations \cite{Ronveaux1995} could illuminate Diophantine properties of exotic constants.

\item \textbf{Extensions.} Natural generalizations include: (a) logarithmic asymptotics when characteristic roots coincide; (b) $q$-analogues where Jacobi theta functions replace Gamma functions.

\item \textbf{Applications.} Black hole quasi-normal modes \cite{Leaver1985} yield three-term recurrences; AGF methods might provide analytical tools if integer slopes appear. Higher-order P-recurrences ($r \geq 3$) could yield linear forms with better Diophantine control.
\end{enumerate}

\begin{thebibliography}{99}

\bibitem{AbramowitzStegun}
M. Abramowitz and I. A. Stegun, \emph{Handbook of mathematical functions}, Dover, 1964.

\bibitem{Barnes1901}
E. W. Barnes, \emph{The theory of the double gamma function}, Philos. Trans. R. Soc. Lond. A \textbf{196} (1901), 265--387.

\bibitem{BirkhoffTrjitzinsky1933}
G. D. Birkhoff and W. J. Trjitzinsky, \emph{Analytic theory of linear difference equations}, Acta Math. \textbf{60} (1933), 1--89.

\bibitem{ChyzakSalvy1998}
F. Chyzak and B. Salvy, \emph{Non-commutative elimination in Ore algebras proves multivariate identities}, J. Symbolic Comput. \textbf{26} (1998), 187--227.

\bibitem{DLMF}
F.~W.~J. Olver, A.~B. Olde Daalhuis, D.~W. Lozier, B.~I. Schneider, R.~F. Boisvert, C.~W. Clark, B.~R. Miller, B.~V. Saunders, H.~S. Cohl, and M.~A. McClain, eds., \emph{NIST Digital Library of Mathematical Functions}, Release 1.2.0 of 2023-12-15, \url{https://dlmf.nist.gov/}.

\bibitem{Finch2003}
S. R. Finch, \emph{Mathematical constants}, Cambridge Univ. Press, 2003.

\bibitem{FlajoletSedgewick2009}
P. Flajolet and R. Sedgewick, \emph{Analytic combinatorics}, Cambridge Univ. Press, 2009.

\bibitem{KauersPaule2011}
M. Kauers and P. Paule, \emph{The concrete tetrahedron: symbolic sums, recurrence equations, generating functions, asymptotic estimates}, Springer, 2011.

\bibitem{Koutschan2010}
C. Koutschan, \emph{Advanced applications of the holonomic systems approach}, PhD thesis, RISC, Johannes Kepler University, Linz, Austria, 2009.

\bibitem{Leaver1985}
E. W. Leaver, \emph{An analytic representation for the quasi-normal modes of Kerr black holes}, Proc. R. Soc. Lond. A \textbf{402} (1985), 285--298.

\bibitem{Loya2017}
P. Loya, \emph{Amazing and aesthetic aspects of analysis}, Springer, 2017.

\bibitem{Norlund1924}
N. E. Nörlund, \emph{Vorlesungen über Differenzenrechnung}, Springer, 1924.

\bibitem{OEIS_A000166}
N. J. A. Sloane, \emph{Sequence A000166 (Subfactorial or rencontres numbers)}, The On-Line Encyclopedia of Integer Sequences, \url{https://oeis.org/A000166}.

\bibitem{OEIS_A006882}
N. J. A. Sloane, \emph{Sequence A006882 (Double factorials n!!)}, The On-Line Encyclopedia of Integer Sequences, \url{https://oeis.org/A006882}.

\bibitem{Ronveaux1995}
A. Ronveaux (ed.), \emph{Heun's differential equations}, Oxford Univ. Press, 1995.

\bibitem{Stanley1999}
R. P. Stanley, \emph{Enumerative combinatorics, volume 2}, Cambridge Univ. Press, 1999.

\bibitem{Temme1996}
N. M. Temme, \emph{Special functions: an introduction to the classical functions of mathematical physics}, Wiley, 1996.

\bibitem{Titchmarsh1939}
E. C. Titchmarsh, \emph{The theory of functions}, 2nd ed., Oxford Univ. Press, 1939.

\bibitem{Wimp1984}
J. Wimp, \emph{Computation with recurrence relations}, Pitman, 1984.

\bibitem{Zeilberger1990}
D. Zeilberger, \emph{A holonomic systems approach to special functions identities}, J. Comput. Appl. Math. \textbf{32} (1990), 321--368.

\bibitem{Zeilberger1991}
D. Zeilberger, \emph{The method of creative telescoping}, J. Symbolic Comput. \textbf{11} (1991), 195--204.

\end{thebibliography}
\end{document}
